\documentclass[]{../../../../tex/classes/homework}
\usepackage{../../../../tex/packages/hebrewSupport}
\usepackage{../../../../tex/packages/mathShortcuts}
\usepackage{../../../../tex/packages/theoremsSupport}


\newcommand\mo {m^*_{\mathrm{open}}}
\newcommand\mc {m^*_{\mathrm{closed}}}

\author{שחר פרץ}
\title{פונקציות ממשיות $\sim$ \textit{תרגיל בית 1}}
\begin{document}
	\maketitle
	\section{}
	נוכיח שבהגדרה של מידה חיצונית ניתן להחליף את התיבות האלמנטריות בתיבות סגורות/פתוחות/אלמנטריות. 
	
	\begin{proof}
		נסמן ב־$m^*$ את המידה שהגדרנו עם תיבות אלמנטריות, ב־$\mo$ לתיבות פתוחות וב־$\mc$ לקבוצות סגורות. תהא $A \subset \R^{d}$. לעת עתה נניח שהיא חסומה ע''י $M$. 
		\begin{itemize}
			\item ראשית נוכיח $m^*(A) = \mc(A)$. משום שתיבות אלמנטריות מוכלות בתוך קבוצות סגורות, ברור ש־$m^*(A) \le \mc(A)$ (מהגדרת האינפימום). נוכיח להוכיח את השוויון מהכיוון השני. יהי $\eg > 0$ ונוכיח ש־$m^*(A) + \eg \ge \mc(A)$. 
			
			יהי כיסוי \textbf{אלמנטרי} $R_j$ של $A$. לכל מלבן אלמנטרי $[a_i, b_i) \times [x_i, y_i)$ נבחין שהקבוצה: 
			\[ E_{r} = \csb{b_i - \frac{r}{2}, b_i + \frac{r}{2}} \times [x_i, y_i] \cup \csb{a_i, b_i} \times \csb{y_i - \frac{r}{2}, y_i + \frac{r}{2}} \]
			מקיימת מתת־אדטיביות ומחסימות ע''י $M$: 
			\[ m^*(E_r) = m(E_r) \le r \cdot (y_i - x_i) + r \cdot (b_i - a_i) \le 2rM \]
			עוד נבחין ש־: 
			\[ R_i = [a_i, b_i) \times [x_i, y_i) \cup E_r \subset [a_i, b_i] \times [x_i, y_i] \doc \ol{R_i} \]
			כאשר $\ol{R_i}$ מלבן סגור. לכל $r > 0$. דהיינו: 
			\[ A = \bigcup_{j = 1}^{\infty}R_j \subset \bigcup_{i = 1}^{\infty}R_j \cup E_{r_j} \subset \bigcup_{i = 1}^{\infty}\ol{R_j} \]
			כאשר $r_j  = \frac{\eg}{2^{n + 2} \cdot M + 1} \ge 0$. מתת־אדטיביות: 
			\[ m\!\cl{\bigcup_{j = 1}^{\infty} \ol{R_j}} \le m\!\cl{\bigcup_{j = 1}^{\infty} R_j} + m\!\cl{\bigcup_{j = 1}^{\infty} E_{r_j}} \]
			ומתת־\ס־אדטיביות: 
			\[ m\!\cl{\bigcup_{j = 1}^{\infty} \ol{R_j}} - m\!\cl{\bigcup_{j = 1}^{\infty} R_j} \le m\!\cl{\bigcup_{j = 1}^{\infty} E_{r_j}} \le \sum_{i = 1}^{\infty}m(E_{r_j}) = \sum_{i = 1}^{n}\frac{2M\eg}{2^{n + 2}M + 1} \le \eg \]
			משום ש־$\inf$ אינווריאנטי להוספה של קבוע $\eg$, סה''כ $m^*(A) + \eg \ge \mc(A)$ לכל $\eg > 0$. נסיק ש־$m^*(A) \ge \mc(A)$ כנדרש. 
			\item עתה נפנה להוכיח ש־$m^*(A) = \mo(A)$. באופן דומה תיבות פתוחות מוכלות בתוך תיבות אלמנטריות ולכן ברור ש־$m^*(A) \ge \mo(A)$ לכל $A \subset \R^{d}$. ניקח כיסוי אלמנטרי $R_j$ של $A$. גם הפעם נגדיר לכל $i \in \N$ את הקבוצה: 
			\[ W_{r} = \csb{a_i - \frac{r}{2}, a_i + \frac{r}{2}} \times [x_i, y_i] \cup \csb{a_i, b_i} \times \csb{x_i - \frac{r}{2}, x_i + \frac{r}{2}} \]
			גם במקרה זה יתקיים $m^*(W_r) \le 2rM$ ונגדיר $r_j$ לאותו הערך. נסמן ב־$R_i\intr$ מלבן פתוח ממדי מלבן $R_i$ (באופן דומה לסעיף הקודם) וגם הפעם $R_j\intr \cup W_j \supset R_j$, מה שגורר באופן דומה:
			\[ m\!\cl{\bigcup_{j = 1}^{\infty} R_j} - m\!\cl{\bigcup_{j = 1}^{\infty} R_j\intr} \le m\!\cl{\bigcup_{j = 1}^{\infty} W_{r_j}} \le \sum_{i = 1}^{\infty}m(W_{r_j}) = \sum_{i = 1}^{n}\frac{2M\eg}{2^{n + 2}M + 1} \le \eg \]
			דהיינו $m^*(A) - \mo(A) \le \eg$ לכל $\eg > 0$ ולכן $m^*(A) \le \mo(A)$ וזה מסיים את ההוכחה. 
		\end{itemize}
	\end{proof}
	
	\section{}
	נוכיח שגרף פונקציה אינטגרבילית רימן הוא ממידה אפס. 
	\begin{proof}
		תהא $f \co [0, 1] \to \R$ אינטגרבילית רימן. מקריטריון דרבו קיימת חלוקה $\Pi$ כך ש־$\daru(f, \Pi) - \dard(f, \Pi) < \eg$. בכל קטע $[x_i, x_{i + 1}]$ בחלוקה נסמן ב־$u_i$ את הערך המקסימלי וב־$d_i$ את הערך המינימלי (מוגדר היטב כי $f$ מקבלת את חסמיה על קבוצה קומפקטית). נסמן את מדד עדינות החלוקה ב־$\lg$ (אורך הקטע המקסימלי בחלוקה), ונבחין שהמלבנים הסגורים $R_i \cod [d_i, u_i] \times [x_i, x_{i +1}]$ הם כיסוי של הגרף של $f$ (ההוכחה פשוטה: לכל $x \in [x_i, x_{i + 1}]$ בהכרח $f(x) \in [d_i, u_i]$ ולכן $(x, f(x)) \in R_i$). איחודם סופי ולכן אלמנטרי. 
		\[ \begin{WithArrows}[groups]
			&m^*(A) \Arrow[i]{הגדרת $\inf$}\\
			\le &\ m^*\!\cl{\bigcup_{i = 1}^{n}R_i} \Arrow[i]{תת־אדטיביות} \\
			\le &\ \sum_{i = 1}^{n}m^*(R_j) \Arrow{מידה חיצונית מזדהה עם מידה אלמנטרית}\\
			\le &\ \sum_{i = 1}^{n}(u_i - d_i) \cdot (x_{i + 1} - x_i) \\
			= &\ \underbrace{\sum_{i = 1}^{n}u_i\Delta x_i}_{\daru(f, \Pi)}  - \underbrace{\sum_{i = 1}^{n}d_i\Delta x_i }_{\dard(f, \Pi)} \le \eg
		\end{WithArrows} \]
		(הערה: למעשה השתמשנו בהגדרה של $m^*$ לפי $\mc$ שהוצגה בשאלה 1, והוכחה שקילות)
		
		סה''כ לכל $\eg > 0$ מתקיים ש־$m^*(A) \le \eg$ ולכן $m^*(A) = 0$. 
	\end{proof}
	
	\section{}
	תהא $A \subset \R$ ממידה אפס. 
	\begin{enumerate}[(A)]
		\item נוכיח ש־$A \times \R \subset \R^{2}$ ממידה אפס. \begin{proof}
			יהי $\eg > 0$. נתבונן בכיסוי $r_i = [a_i, b_i)$ כיסוי של $A$ כך ש־$\sum m(r_i) < \min\{\frac{\eg}{2}, 1\}$ (קיים כזה מהגדרת $\inf$). נגדיר: 
			\[ R_i = r_i \times [-\eg_i, \eg_i) \]
			כאשר $\eg_n = \frac{\min\{\eg, 1\}}{2^{n + 4}} > 0$. נבחין ש־$m^*(R_i) = m(R_i) = 2\eg_nm(r_i)$. עתה מסאב־\ס־אדטיביות: 
			\[ m^*(A) \le m\!\cl{\bigcup_{i = 1}^{\infty}R_i} \le \sum_{i = 1}^{\infty}2 \eg_i m( r_i) = \cdots \]
			מא''ש הממוצעים, לכל $a, b \in \R_{\ge 0}$: 
			\[ \frac{a + b}{2} \ge \sqrt{ab} \implies (a + b)^{2} \ge 4ab \implies a^2 \!+ b^2 \ge 3ab \]
			מהתוצאה הזו, ומכיוון שבגלל ש־$m(\bigcup r_i) < 1\}$ אז גם $m(r_i) < 1$, ולכן $m(r_i)^{2} \le m(r_i)$ וכנ''ל בעבור $\eg_i$: 
			\[ \cdots \le \frac{2}{3}\sum_{i = 1}^{\infty}\cl{\eg_i^2  + m(r_i)^2} \le \sum_{i = 1}^{\infty}\eg_i^{2} + \sum_{i = 1}^{\infty}m(r_i)^{2} \le \underbrace{\sum_{i = 1}^{\infty}\eg_i}_{\mathclap{\le \sum\frac{\eg}{2^{n + 4}} \le \frac{\eg}{2}}} + \underbrace{\sum_{i = 1}^{n}m(r_i)}_{\le \frac{\eg}{2}} \le \frac{\eg}{2} = \eg \]
			כלומר לכל $\eg > 0$ מתקיים $m^*(A) \le \eg$ כנדרש. 
		\end{proof}
		\item תהא $f \co \R \to \R$ פונקציה $L$־לפישיץ. נוכיח ש־$f(A)$ ממידה אפס. \begin{proof}
			יהי $\eg> 0$. ניקח כיסוי $r_i =[a_i, b_i)$ של $A$ כך ש־$\sum m(r_i) \le \frac{\eg}{L}$ שקיים מהיותה זניחה. נגדיר $R_i = f(r_i)$. נבחין ש־: 
			\[ f(A) \subset f\!\cl{\bigcup_{i = 1}^{\infty} r_i} = \bigcup_{i = 1}^{\infty}f(r_i) = \bigcup_{i = 1}^{\infty}R_i \]
			לכן $R_i$ כיסוי של $f(A)$. בגלל שאם $r_i = [a_i, b_i)$ אז $\sof{f(b_i) - f(a_i)} \le L\sof{b_i - a_i}$ מהגדרת $m$ למלבנים $m(f(r_i)) \le Lm(r_i)$. מתת־\ס־אדטיביות ומשום ש־$m^*$ מזדהה עם $m$ בעבור אלמנטריים, נקבל ש־: 
			\[ \sum_{i = 1}^{\infty}m^*(R_i) = \sum_{i = 1}^{\infty}m(f(r_i)) \le L\sum_{i = 1}^{\infty}m(r_i) \le L \cdot \frac{\eg}{L} = \eg \]
			מהגדרת אינפימום ומהגדרת מידה, $f(A)$ ממידה אפס. 
		\end{proof}
	\end{enumerate}
	
	\section{}
	תהא $\cc$ קבוצת קנטור. 
	\begin{enumerate}[(A)]
		\item נראה ש־$\cc$ היא nowhere dense. נניח בשלילה שישנה פתוחה $(a, b)$ כך ש־$\Cl(\cc \cap (a, b)) = (a, b)$. ניעזר ב־$C_i$ באופן שבו הוגדרו בתרגול. נסמן $m = b - a$. 
		קל לראות מהבנייה המפורשת מהתרגול שהקבוצה $(a, b)$ הגדולה ביותר שמוכלת ב־$C_{i + 2}$ היא מהגודל $\frac{1}{3^{i}}$. לכן עבור $N = \log_3(\sof{b - a})$, נקבל שישנו קטע פתוח $(a', b') \subset (a, b)$ כך ש־$(a', b') \cap C_{i + 2} = \vn$, ולכן אין סדרה ב־$C_{i + 2}$ שמתכנסת ל־$\frac{b' + a'}{2}$. לכן $C_{i + 2}$ לא צפופה ב־$(a, b)$, ולכן גם $\cc \subset C_{i + 2}$ לא צפופה ב־$(a, b)$. 
		\item נוכיח שאין ב־$\cc$ נקודות מבודדות. \begin{proof}
			אם $a$ קצה קטע ב־$C_n$, בניית $C_{n + 1}$ הוגדרה ככזו בה לוקחים את השליש האמצעי בכל קטע כזה, ולכן ב־$C_{n + 1}$ היא לא תוסר (שכן היא בקצוות ולא בשליש האמצעי). יתרה מכך, היא תיוותר קצה קטע (הקטע יחולק ל־2 חלקים שווים). לכן באינדוקציה על $n$ היא שייכת ל־$C_m$ לכל $m \in \N$ ומהגדרת חיתוך גם ל־$\cc$. 
			
			עתה, נניח בשלילה שקיימת נקודה מבודדת $x$, כלומר ישנה $(a, b)$ פתוחה כך ש־$(a, b) \cap \cc = \{x\}$. מהגדרה $x \in C_n$ לכל $n$. נבחין שמשום שב־$C_n$ הקטעים באורך $3^{-n}$ (פונ' ששואפת ל־$0$), ישנו $N$ גדול דיו כך שאורך הקטעים ב־$C_n$ קטן מ־$\frac{b - a}{4}$. בפרט $x \in C_N$, ומהגדרת הקבוצה $C_N$, היא נמצאת בקטע כלשהו באחד מ־$2^{N}$ הקטעים בקבוצה. קצה קטע יסומן ב־$a$ (בה''כ $a \neq x$, אחרת נוכל לבחור את קצה הקטע השני) והוא נמצא ב־$\cc$ (הוכחנו זאת בפסקה הקודמת) ובחרנו את $N$ כך ש־$\sof{x - a} < b - a$, דהיינו $a \in (a, b) \cap \cc$ למרות ש־$a \neq x$ וזו סתירה. 
		\end{proof}
	\end{enumerate}
	
	\section{}
	נוכיח שהייצוג הטרנארי קיים ויחיד עם מקדמים $0, 2$ אמ''מ $x \in \cc$, ונסיק שקבוצת קנטור מעוצמת הרצף. 
	\begin{proof}
		למה: בהינתן מספרים $x_1 \dots x_n$ כך ש־$x_1 \dots x_{n n} \in \{0, 2\}$, המספר $x^{(n)} \cod \sum_{k = 1}^{\infty}\frac{x_k}{3^{k}}$ נמצא בקבוצה ה־$x_1,x_2,x_3 \dots x_n$ בייצוג בינארי ב־$C_N$, כאשר $2$־מים נתפשים כ־$1$־ים. נוכיח את הלמה באינדוקציה. 
		
		נתבונן ב־$x^{(n -1)}$, ונניח שהוא נמצא במיקום ה־$N$. נבחין שב־$C_n$ יש $2N$ מספרים לפני שמגיעים לקטע שלו (ולכן ה־left binary shift). עתה אם $x_n = 0$ הוא יהיה בחצי הראשון של קטע זה, ולכן במקום ה־$2N$, אחרת במיקום ה־$2N + 1$ (במקרה ש־$x_n = 2$). בהתאם לכך ספרת האחדות בייצוג הבינארי, כנדרש. 
		
		עתה, נראה את הגרירה הדו־כיוונית הנדרשת. 
		\begin{itemize}
			\item[$\implies$]נניח ש־$x \in \cc$. נבחין שמשום שלכל $n$ יש מספר קטע ב־$C_N$ אליו $x$ שייך, הלמה שהוכחנו מגדירה סדרה של מספרים $x_k$ כך ש־$x^{(n)}$ נמצא באותו הקטע ב־$C_n$ כמו $x$. נבחין ש־$x^{(n)} \to x$ שכן $x^{(n)}$ ו־$x$ נמצאים יחדיו ב־$C_n$ בקטע שאורכו $3^{-n}$ ולכן $\sof{x^{(n)} - x} \le 3^{-n} \to 0$. לכן הטור מתכנס ל־$x$, כלומר יש לו ייצוג טרנארי מתאים. 
			\item[$\impliedby$]יהי $x$ בעל ייצוג טרנארי עם ספרה $x_k \in \{0, 2\}$. נוכיח ש־$x \in \cc$. נבחין ש־$x$ הוא גבול של מספרים $x^{(n)}$ כך ש־$x^{(n)} \in C_n$ קצוות קטעים ב־$C_n$, והראנו כחלק מ־4(ב) שקצוות קטעים גם ב־$\cc$ לכן $x^{(n)} \in \cc$. לכן $x$ גבול של איברים ב־$\cc$ ומשום ש־$\cc$ סגורה, $x \in \cc$ כנדרש. 
		\end{itemize}
		
		סה''כ יש התאמה חח''ע ועל בין ייצוג טרנארי $\{0, 2\}^{\N}$ ל־$\cc$, לכן: 
		\[ \sof{\cc} = \sof{\{0, 2\}^{\N}} = 2^{\az} = \con \]
		כנדרש. 
	\end{proof}
	
	\section{}
	נוכיח שכל $A \subset \R$ הוא איחוד לכל היותר בן־מנייה של קטעים פתוחים זרים. 
	
	\textbf{הערה: }משפט זה הוכח בקלטות של 2016, שבהן צפיתי וסיכמתי. אז אני מצרף את ההוכחה מהסיכום. 
	
	\begin{proof}
		נתונה קבוצה $A$ פתוחה. נביט באוסף הקטעים המקסימליים ביחס ל־$A$. 
		\begin{itemize}
			\item \textbf{הקטעים הנ''ל פתוחים: }טרוויאלי, קטע פתוח הוא קבוצה פתוחה. 
			\item \textbf{איחודם מכסה את $\bm{A}$: }הראנו בהרצאה שסביב כל נקודה קיים קטע $I_x$ פתוח מקסימלי שמוכל ב־$A$ (שנדרשת להיות פתוחה). 
			\item \textbf{הקטעים הנ''ל זרים: }נניח ש־$I, J \subset A$ שני קטעים מקסימליים שנחתכים, ובשלילה $I \cap J \neq \vn$. אזי $I \cup J \subset A$ פתוח, ומכיל את הקטע המקסימלי $I$ כלומר $I = I \cup J$. באותו האופן $J = I \cup J = I$ ולכן הם שווים. 
			\item \textbf{האיחוד לכל היותר בן בנייה: }כי $\R$ הוא CII – בכל קטע ישנו רציונלי, הרציונלים שווים זה מזה (הם זרים) ומאקסיומת הבחירה הבת מנייה הגדרנו העתקה חח''ע לקבוצה בת מנייה. 
		\end{itemize}
	\end{proof}
	
	\ndoc
\end{document}